\section{Overview of Quantum Mechanics}

A pure state of a qubit can be represented by a 2-dimensional ket vector 
$\ket{\psi} \in \mathbb{C}^{2}$:
\[
\ket{\psi} = \alpha \ket{0} + \beta \ket{1},
\]
where $\alpha, \beta \in \mathbb{C}$ satisfy
\[
|\alpha|^2 + |\beta|^2 = 1.
\]

Whenever this condition on $\alpha$ and $\beta$ is satisfied, we say that the state vector is \emph{normalized}. The complex numbers $\alpha$ and $\beta$ are called the \emph{probability amplitudes} of $\ket{\psi}$.

\begin{theorem}[General Pure State]
For every $\theta \in \mathbb{R}$, the state
\[
\ket{\psi} = \sin\theta \ket{0} + \cos\theta \ket{1}
\]
is a valid pure state.
\end{theorem}

\begin{theo}[Geneal Pure State]{General Pure state}
We verify the normalization condition:
\[
|\sin\theta|^2 + |\cos\theta|^2
= \sin^2\theta + \cos^2\theta
= 1.
\]
Hence, the state is normalized and therefore represents a valid pure state.
\end{theo}

More generally, quantum information can be encoded in higher-dimensional systems. A $d$-dimensional quantum system is called a \emph{qudit}. Its pure state can be written as
\[
\ket{\psi} = \sum_{i=0}^{d-1} \alpha_i \ket{i},
\]
where $\alpha_i \in \mathbb{C}$ satisfy the normalization condition
\[
\sum_{i=0}^{d-1} |\alpha_i|^2 = 1.
\]

\begin{definition}[Standard Basis]
The \emph{standard basis}, also known as the \emph{computational basis}, of 
$\mathbb{C}^2$ is the orthonormal basis
\[
\mathcal{S} = \{\ket{0}, \ket{1}\},
\]
where
\[
\ket{0} =
\begin{pmatrix}
1\\
0
\end{pmatrix},
\qquad
\ket{1} =
\begin{pmatrix}
0\\
1
\end{pmatrix}.
\]
\end{definition}

There are many other bases for $\mathbb{C}^2$. Another important basis is the Hadamard basis.

\begin{definition}[Hadamard Basis]
The \emph{Hadamard basis} of $\mathbb{C}^2$ is the orthonormal basis
\[
\mathcal{H} = \{\ket{+}, \ket{-}\},
\]
where
\[
\ket{+}
=
\frac{1}{\sqrt{2}}(\ket{0}+\ket{1})
=
\frac{1}{\sqrt{2}}
\begin{pmatrix}
1\\
1
\end{pmatrix},
\]
and
\[
\ket{-}
=
\frac{1}{\sqrt{2}}(\ket{0}-\ket{1})
=
\frac{1}{\sqrt{2}}
\begin{pmatrix}
1\\
-1
\end{pmatrix}.
\]
\end{definition}

We now verify that this is indeed an orthonormal basis using bra-ket notation:
\[
\braket{+}
=
\frac{1}{2}
\begin{pmatrix}
1 & 1
\end{pmatrix}
\begin{pmatrix}
1\\
1
\end{pmatrix}
=
\frac{1}{2}(2)
=
1.
\]

Similarly,
\[
\braket{-}
=
\frac{1}{2}
\begin{pmatrix}
1 & -1
\end{pmatrix}
\begin{pmatrix}
1\\
-1
\end{pmatrix}
=
1,
\]
and
\[
\braket{+}{-}
=
\frac{1}{2}
\begin{pmatrix}
1 & 1
\end{pmatrix}
\begin{pmatrix}
1\\
-1
\end{pmatrix}
=
0.
\]

Hence, $\mathcal{H}$ is an orthonormal basis of $\mathbb{C}^2$.